Emmy kan niet langer zich in haar geest verliezen De migraine brenmgt haar naar een vorig leven om van lijdzaam ondergaan naar dader te gaan die ze weer op een meer afstand beschouwd dan zich onbewust te verliezen zondeer nog reflectie fo beschouwing of contact met de wereld meer. Door Martino krijgt ze van de filosofie in haar werelds deel door de direct ervaring toegang tot een dieper niveau en dat brengt haar bij de ziel die alleen maar zijn is/beleven en niet meer denken of wereldse invloeden.

je beschrijft een boeiende paradox, namelijk dat trauma, ondanks de pijn, beweging kan brengen in patronen die anders mogelijk verstard zouden blijven
alsof het trauma een scharnier wordt in de tijd, waarmee plotseling het ogenschijnlijk onwrikbare zijn beweert te kunnen verschuiven

wanneer een vastgeroest gedragspatroon — bijvoorbeeld chronische terughoudendheid — door de schok van het trauma wordt opgeschrikt, kan dat soms leiden tot een compenserende overgave aan het tegenovergestelde (zoals plotselinge openheid of zelfs impulsiviteit)
dit zie je ook in sommige therapeutische modellen terug, waar de abrupte verstoring door trauma oude defensies ontregelt en iemand noodgedwongen alternatieven uitprobeert

vergelijk het met een deur die eindeloos op slot zat — het trauma forceert een opening
soms met geweld, soms met precisie
en dan blijkt er iets mogelijk wat voorheen ondenkbaar leek
maar de angst voor het onbekende achter die deur blijft meespelen

tegelijkertijd draait het niet alleen om de verschuiving van het gedrag zelf, maar ook om de integratie ervan
het ‘nut’ van trauma ontstaat pas werkelijk als iemand het nieuwe terrein niet ziet als “enkel compensatie”, maar als een kans tot hernieuwde vrijheid
het is de transformatie, niet de oppervlakkige omkering, waarin de kracht van de ervaring schuilt

en toch: het openen van die deur leidt niet noodzakelijkerwijs tot verlichting of genezing
soms komt er enkel een koude tocht binnen
de kunst is wellicht het openen te begeleiden,
zodat het nieuwe niet weer stolt tot een nieuw pantser

wil je deze reflectie verwerken als commentaar bij een canto, of als onderdeel van een bespiegeling in Emmy Nativus?
of zoek je een verdere uitwerking in de vorm van een essayistische passage of een dramatische dialoog?

Vanuit Emmy’s gêne kwam weer de migraine. De regressiescène toonde warme beelden. Hoe ze samen speelden en momenten deelden met haar zuster Emma. Het verzachtte bijna
het pijnlijke drama. De tijd bleef teruggaan en kwam bij de kiem aan van haar aardse bestaan. Op naar het paradijs ging toen haar zielenreis en hopte sprongsgewijs om zich te begeven naar haar tweede leven om her te beleven.

Emmy bevond zichzelf al neukend op een schelf in 1411 uit politiek gewin met een jonge gravin, echter tegen diens zin. Hij moest terreur zaaien door adel te naaien om de macht te graaien voor zijn baas Braccio, een condottiero berucht in Arezzo, de stad waar hij verbleef. Hij, die Emmy omschreef als haar altergeilneef, was ooit een arme boer. Hij werd een drinkebroer en koos de rovers toer in de streek Toscane waar hij inhumane, geenszins toegestane, misdaden ging plegen. In aanzien gestegen kruisten elkaars wegen de wilde Martino, als haar alter ego, en krijgsman Braccio.

Martino was beducht en al lang op de vlucht na een strafbare klucht. Hij greep in een bordeel in het snoodste stadsdeel de pooier bij de keel toen hij ervan baalde dat hij goed betaalde, maar dat niet vertaalde in een passende hoer. Dus vermoordde de boer de man in een krachttoer en ging de werkkrachten een voor een verkrachten. Vaak knepen stadswachten een oogje voor hem dicht. Het kwaad hem toegedicht  had voldoende gewicht, maar met deze misdaad werden ze wel kordaat. De dood werd de schuldmaat.

De huurlingenleider zag in de zwartrijder een begaafde strijder en lijfde hem win-win als privé schildknaap in. Zijn vadertje Stalin regelde zijn vrijheid in ruil voor zijn trouwheid tot zijn dood in de strijd. Bungelend aan een touw of een lonende trouw? Tis wel klaar wat hij wou. De kunst van het vechten leren was voor knechten een van de voorrechten. Hij verzorgde zijn paard, droeg zijn schild en zijn zwaard, maar het meeste was waard de brute ervaring. De wrede veldtraining en scherpe waarneming.

Niet bang om te doden of martelmethoden, die waren verboden, bezield uit te voeren, wist de tong te roeren. Hij brak de contouren van zijnde een dienaar bij elk extreem gevaar en werd tot slot zowaar een volwaardig soldaat. Braccio koos het kwaad. Hij deed het adequaat. Zoals het verkrachten om daarmee te trachten alle tegenkrachten te demotiveren en te penetreren. Een lesje te leren.

Zo lag hij in het hooi met een gravin bloedmooi en dumpte zijn rotzooi. Emmy kon het douwen niet langer aanschouwen. Haar blik ging vernauwen en de vortex kwam weer. De aangereikte leer ging over “lust” dit keer.

Door het kijken naar Martino beseft Emmy dat zij ooit zo ook geweest heeft gaan voor genot, voor lust. En herinnert ze zich ook het gevoel wat ze al eerder voor het eerst heeft ervaren bij de jager. En dat dat een hele sterke kracht is. En eigenlijk heeft ze ook, ondanks dat alle oordelen er zijn, het besef dat zij ooit Martino was, betekent dat ze dat niet van zich weg kan gooien. En daardoor gaat ze terug naar het gevoel toen ze Martino in het ogen keek toen hij zijn hoogtepunt bereikte, dacht ze terug aan haar eerdere hoogtepunten die ze voor het eerst bij de jager had. En daardoor viel het oordelen weg en begreep ze zijn drive zonder dat goed te keuren of af te keuren. Maar ze leverde zich in en zag ook de ogen van de jager erin terug. En dus werd het niet een persoonlijk iets meer, maar over het algemeen de sensatie van de climax. En dat bracht haar in een hogere zweer, waarbij het niet meer ging om mensen of daden of oordelen, maar puur om de sensatie van de climax.

De climax voelt ze ineens met hem mee alsof ze in hem verdwijnt en dan brengt haar de hemel
